% Options for packages loaded elsewhere
\PassOptionsToPackage{unicode}{hyperref}
\PassOptionsToPackage{hyphens}{url}
\PassOptionsToPackage{dvipsnames,svgnames,x11names}{xcolor}
%
\documentclass[
  letterpaper,
  DIV=11,
  numbers=noendperiod]{scrartcl}

\usepackage{amsmath,amssymb}
\usepackage{iftex}
\ifPDFTeX
  \usepackage[T1]{fontenc}
  \usepackage[utf8]{inputenc}
  \usepackage{textcomp} % provide euro and other symbols
\else % if luatex or xetex
  \usepackage{unicode-math}
  \defaultfontfeatures{Scale=MatchLowercase}
  \defaultfontfeatures[\rmfamily]{Ligatures=TeX,Scale=1}
\fi
\usepackage{lmodern}
\ifPDFTeX\else  
    % xetex/luatex font selection
\fi
% Use upquote if available, for straight quotes in verbatim environments
\IfFileExists{upquote.sty}{\usepackage{upquote}}{}
\IfFileExists{microtype.sty}{% use microtype if available
  \usepackage[]{microtype}
  \UseMicrotypeSet[protrusion]{basicmath} % disable protrusion for tt fonts
}{}
\makeatletter
\@ifundefined{KOMAClassName}{% if non-KOMA class
  \IfFileExists{parskip.sty}{%
    \usepackage{parskip}
  }{% else
    \setlength{\parindent}{0pt}
    \setlength{\parskip}{6pt plus 2pt minus 1pt}}
}{% if KOMA class
  \KOMAoptions{parskip=half}}
\makeatother
\usepackage{xcolor}
\usepackage{svg}
\setlength{\emergencystretch}{3em} % prevent overfull lines
\setcounter{secnumdepth}{-\maxdimen} % remove section numbering
% Make \paragraph and \subparagraph free-standing
\makeatletter
\ifx\paragraph\undefined\else
  \let\oldparagraph\paragraph
  \renewcommand{\paragraph}{
    \@ifstar
      \xxxParagraphStar
      \xxxParagraphNoStar
  }
  \newcommand{\xxxParagraphStar}[1]{\oldparagraph*{#1}\mbox{}}
  \newcommand{\xxxParagraphNoStar}[1]{\oldparagraph{#1}\mbox{}}
\fi
\ifx\subparagraph\undefined\else
  \let\oldsubparagraph\subparagraph
  \renewcommand{\subparagraph}{
    \@ifstar
      \xxxSubParagraphStar
      \xxxSubParagraphNoStar
  }
  \newcommand{\xxxSubParagraphStar}[1]{\oldsubparagraph*{#1}\mbox{}}
  \newcommand{\xxxSubParagraphNoStar}[1]{\oldsubparagraph{#1}\mbox{}}
\fi
\makeatother


\providecommand{\tightlist}{%
  \setlength{\itemsep}{0pt}\setlength{\parskip}{0pt}}\usepackage{longtable,booktabs,array}
\usepackage{calc} % for calculating minipage widths
% Correct order of tables after \paragraph or \subparagraph
\usepackage{etoolbox}
\makeatletter
\patchcmd\longtable{\par}{\if@noskipsec\mbox{}\fi\par}{}{}
\makeatother
% Allow footnotes in longtable head/foot
\IfFileExists{footnotehyper.sty}{\usepackage{footnotehyper}}{\usepackage{footnote}}
\makesavenoteenv{longtable}
\usepackage{graphicx}
\makeatletter
\def\maxwidth{\ifdim\Gin@nat@width>\linewidth\linewidth\else\Gin@nat@width\fi}
\def\maxheight{\ifdim\Gin@nat@height>\textheight\textheight\else\Gin@nat@height\fi}
\makeatother
% Scale images if necessary, so that they will not overflow the page
% margins by default, and it is still possible to overwrite the defaults
% using explicit options in \includegraphics[width, height, ...]{}
\setkeys{Gin}{width=\maxwidth,height=\maxheight,keepaspectratio}
% Set default figure placement to htbp
\makeatletter
\def\fps@figure{htbp}
\makeatother

\KOMAoption{captions}{tableheading}
\makeatletter
\@ifpackageloaded{caption}{}{\usepackage{caption}}
\AtBeginDocument{%
\ifdefined\contentsname
  \renewcommand*\contentsname{Table of contents}
\else
  \newcommand\contentsname{Table of contents}
\fi
\ifdefined\listfigurename
  \renewcommand*\listfigurename{List of Figures}
\else
  \newcommand\listfigurename{List of Figures}
\fi
\ifdefined\listtablename
  \renewcommand*\listtablename{List of Tables}
\else
  \newcommand\listtablename{List of Tables}
\fi
\ifdefined\figurename
  \renewcommand*\figurename{Figure}
\else
  \newcommand\figurename{Figure}
\fi
\ifdefined\tablename
  \renewcommand*\tablename{Table}
\else
  \newcommand\tablename{Table}
\fi
}
\@ifpackageloaded{float}{}{\usepackage{float}}
\floatstyle{ruled}
\@ifundefined{c@chapter}{\newfloat{codelisting}{h}{lop}}{\newfloat{codelisting}{h}{lop}[chapter]}
\floatname{codelisting}{Listing}
\newcommand*\listoflistings{\listof{codelisting}{List of Listings}}
\makeatother
\makeatletter
\makeatother
\makeatletter
\@ifpackageloaded{caption}{}{\usepackage{caption}}
\@ifpackageloaded{subcaption}{}{\usepackage{subcaption}}
\makeatother

\ifLuaTeX
  \usepackage{selnolig}  % disable illegal ligatures
\fi
\usepackage{bookmark}

\IfFileExists{xurl.sty}{\usepackage{xurl}}{} % add URL line breaks if available
\urlstyle{same} % disable monospaced font for URLs
\hypersetup{
  pdftitle={Economía Política},
  colorlinks=true,
  linkcolor={blue},
  filecolor={Maroon},
  citecolor={Blue},
  urlcolor={Blue},
  pdfcreator={LaTeX via pandoc}}


\title{Economía Política}
\usepackage{etoolbox}
\makeatletter
\providecommand{\subtitle}[1]{% add subtitle to \maketitle
  \apptocmd{\@title}{\par {\large #1 \par}}{}{}
}
\makeatother
\subtitle{Unidad 5. Estos son mis principios. Si no le gu\$tan, tengo
otros. Influencia, dinero y grupos de interés}
\author{}
\date{}

\begin{document}
\maketitle

\renewcommand*\contentsname{Table of contents}
{
\hypersetup{linkcolor=}
\setcounter{tocdepth}{2}
\tableofcontents
}

\begin{quote}
If money go before, all ways do lie open.\\
\textbf{{[}William Shakespeare, Ford, The Merry Wives of Windsor, Act 2
Scene 2 (1602){]}}
\end{quote}

\section{\texorpdfstring{\textbf{Financiamiento de campañas y
lobbying}}{Financiamiento de campañas y lobbying}}\label{financiamiento-de-campauxf1as-y-lobbying}

\begin{itemize}
\tightlist
\item
  Actividad pre-electoral y post-electoral
\item
  Regulación versus desregulación
\item
  ¿Comprando la política o comprando influencia?
\end{itemize}

\subsection{Financiamiento de
campañas}\label{financiamiento-de-campauxf1as}

\begin{itemize}
\tightlist
\item
  ¿Y por casa cómo andamos? En Argentina, las dos leyes más importantes
  que regulan la actividad electoral y el financiamiento de partidos
  politicos son la Ley 26571 (2009) --hito principal el establecimiento
  de las PASO- y la Ley 26215 (2007) y sus modificatorias
\item
  ¿El \emph{lobbying} está institucionalizado en Argentina? No!

  \begin{itemize}
  \tightlist
  \item
    hay algunas iniciativas pero no hay una Ley Nacional
  \item
    Régimen de Acceso a la Información Pública (2003, luego Ley 27275
    (2016))

    \begin{itemize}
    \tightlist
    \item
      Link: \href{https://audiencias.mininterior.gob.ar/}{Registro Único
      de Audiencias de Gestión de Intereses}
    \end{itemize}
  \end{itemize}
\end{itemize}

\subsection{Financiamiento de campañas
(cont.)}\label{financiamiento-de-campauxf1as-cont.}

\begin{itemize}
\tightlist
\item
  Importante notar la deficiencia en la sistematización, publicación y
  difusíon de estadísticas en relación a la actividad de financiamiento
  partidario
\item
  Sin embargo, algunos avances en años recientes Link:
  \href{https://www.electoral.gov.ar/financiamiento/aportes-privados.php}{Financiamiento
  de Campañas - Aportes privados}
\item
  En Argentina pueden aportar personas físicas y jurídicas --ambas con
  limitaciones. No pueden aportar personas físicas/jurídicas extranjeras
  y tampoco gobiernos y entidades públicas extranjeras
\end{itemize}

\subsection{Financiamiento de campañas
(cont.)}\label{financiamiento-de-campauxf1as-cont.-1}

\begin{figure}

{\centering \includesvg{fig/fig-05-007.svg}

}

\caption{Contribuciones por sectores de actividad - Personas jurídicas}

\end{figure}%

\subsection{Gasto en publicidad}\label{gasto-en-publicidad}

\begin{itemize}
\tightlist
\item
  El gasto en \textbf{publicidad tradicional} --principalmente radio y
  TV- puede consultarse en sitios que agregan y monitorean esta
  información: \href{https://www.opensecrets.org/political-ads}{Gasto en
  publicidad tradicional}
\item
  Interesante mirar también la actividad de donación y contribuciones en
  forma de \textbf{publicidad digital} --gasto en publicidad digital
  tomando con fuente Google y Facebook
\item
  Recurso útil:
  \href{https://www.opensecrets.org/online-ads/map/google}{Gasto en
  publicidad digital}
\item
  Interesante correlacionar el gasto en publicidad digital con ultimo
  resultado electoral/encuestas de intención de voto
\end{itemize}

\subsection{\texorpdfstring{El rol de la \emph{dark
money}}{El rol de la dark money}}\label{el-rol-de-la-dark-money}

\begin{itemize}
\tightlist
\item
  En lingo especialista, \textbf{dark money} refiere al gasto con
  intención de afectar los resultados políticos en el cual la
  fuente/origen de los fondos no es revelado
\item
  En EEUU esto es posible por la existencia de ONGs políticamente
  activas (tipología 501(c)(4)) que no tienen obligacion de declarar sus
  donantes
\item
  Esto ``explotó'' luego de 2010 de producido el fallo de la Corte
  ``Citizens United v. FEC Supreme Court'' -- que basicamente permitió
  la existencia y proliferació de estas grupos políticamente activos
\end{itemize}

\subsection{\texorpdfstring{El rol de la \emph{dark money}
(cont.)}{El rol de la dark money (cont.)}}\label{el-rol-de-la-dark-money-cont.}

\begin{quote}
\textbf{Citizens United v. FEC Supreme Court.} En 2007, en FEC v.
Wisconsing Right to Life la Corte Suprema falló que las organizaciones
sin fines de lucro podía usar fondos para publicidad antes de las
elecciones siempre y cuando estas publicidades no se expresaran
explícitamente a favor (o en contra) de algún candidato. En el año 2010,
la Corte Suprema falló que la prohibición que pesaba sobre
contribuciones politicas independientes violaba el derecho a la libre
expresión contenido en la Primera Enmienda. Este fallo permitió
deliberadamente el uso de fondos para publicidad que se expresara
explícitamente a favor (o en contra) de algún candidato. Desde este
fallo, se produjo la década de mayor gasto en campañas en la historia de
los EEUU.
\end{quote}

\subsection{\texorpdfstring{El rol de la \emph{dark money}
(cont.)}{El rol de la dark money (cont.)}}\label{el-rol-de-la-dark-money-cont.-1}

\begin{figure}

{\centering \includegraphics{fig/fig-05-006.png}

}

\caption{La explosión de gasto de los Super PACs seguida del fallo
Citizens United v. FEC}

\end{figure}%

\subsection{\texorpdfstring{El rol de la \emph{dark money}
(cont.)}{El rol de la dark money (cont.)}}\label{el-rol-de-la-dark-money-cont.-2}

\begin{itemize}
\tightlist
\item
  ¿Quiénes son algunos de los más importantes grupos de \emph{dark
  money}?

  \begin{itemize}
  \tightlist
  \item
    La Cámara de Comercio de EEUU; Crossroads GPS (Karl Rove); Americans
    for Prosperity (Koch Brothers)
  \end{itemize}
\item
  Un tema adicional es la influencia del \textbf{dinero extranjero} en
  las actividades de estos grupos --legalmente está prohibido donaciones
  de personas físicas y jurídicas extranjeras:

  \begin{itemize}
  \tightlist
  \item
    La Cámara de Comercio de EEUU recibe dinero de membresía de empresas
    que representa --muchas incluso de propiedad extranjera
  \item
    La National Rifle Asociation, alineado con el partido republicano,
    recibe cuantiosas sumas de fabricantes extranjeros como SIG Sauer y
    Beretta
  \item
    El American Chemistry Council también tiene a muchas empresas
    extranjeras entre sus miembros
  \end{itemize}
\end{itemize}

\subsection{\texorpdfstring{El \emph{lobbying} como
actividad}{El lobbying como actividad}}\label{el-lobbying-como-actividad}

\begin{itemize}
\tightlist
\item
  El \emph{lobbying} en EEUU está implicitamente (originalmente)
  aceptado a partir de la Primera Enmienda pero adoptó status legal en
  1946 con la Federal Regulation of Lobbying Act que fue luego
  reemplazada por la Lobbying Disclosure Act de 1995
\item
  La High Authority for Transparency in Public Life ha publicado un
  estudio para 41 países documentando las provisiones en relación a la
  actividad de \emph{lobbying}

  \begin{itemize}
  \tightlist
  \item
    Sólo 17 países tienen una ley nacional de lobby (regulación de
    intereses especiales)
  \item
    22 países no tienen ley nacional de lobby
  \item
    El resto se rigen por leyes local o adhesión a leyes supranacionales
  \end{itemize}
\end{itemize}

\subsection{\texorpdfstring{El \emph{lobbying} como actividad
(cont.)}{El lobbying como actividad (cont.)}}\label{el-lobbying-como-actividad-cont.}

\begin{figure}

{\centering \includegraphics{fig/fig-05-002.jpg}

}

\caption{Composición del gastos en financiamiento de campañas y en
lobbying por grupos etarios (generaciones)}

\end{figure}%

\section{\texorpdfstring{\textbf{Modelos básicos de competencia
electoral con grupos de
interés}}{Modelos básicos de competencia electoral con grupos de interés}}\label{modelos-buxe1sicos-de-competencia-electoral-con-grupos-de-interuxe9s}

\subsection{Grupos de presión y poder
político}\label{grupos-de-presiuxf3n-y-poder-poluxedtico}

\begin{itemize}
\tightlist
\item
  Ahora suponga que algunos grupos de individuos en la sociedad tienen
  capacidad de organizarse como \textbf{grupos de presión} (lobbies) y
  aportar contribuciones de campaña
\item
  Todos los grupos tienen la misma densidad ideológica \(\phi{j}=\phi\)
  --el poder no proviene de su sensibilidad ideológica
\item
  Variable \(O^{j}\) igual a 1 si grupo se organiza como lobby; 0 en
  caso contrario

  \begin{itemize}
  \tightlist
  \item
    cada individuo de un lobby contribuye \(C^{j}_{P}\) recursos al
    candidato \(P\)
  \end{itemize}
\end{itemize}

\subsection{Grupos de presión y poder político
(cont.)}\label{grupos-de-presiuxf3n-y-poder-poluxedtico-cont.}

\begin{itemize}
\tightlist
\item
  Cantidad total de fondos recibida por \(P\):
\end{itemize}

\begin{equation}
C_{P}=\sum_{j=1}^{n} O^{j}\alpha^{j}C^{j}_{P}
\end{equation}

\begin{itemize}
\tightlist
\item
  Suponemos que candidato recibe contribuciones \emph{luego} de anunciar
  su plataforma de política pero \emph{antes} que la preferencia general
  \(\delta\) sea conocida
\end{itemize}

\subsection{Grupos de presión y poder político
(cont.)}\label{grupos-de-presiuxf3n-y-poder-poluxedtico-cont.-1}

\begin{itemize}
\tightlist
\item
  Pero ahora \(\delta\) tiene dos componentes:

  \begin{itemize}
  \tightlist
  \item
    \(\tilde{\delta}\) igual al modelo anterior
  \item
    componente \(h(C_{B}-C_{A})\) --introduce el hecho de que las
    contribuciones permiten mejorar popularidad general con un parametro
    de efectividad \(h\)
  \end{itemize}
\end{itemize}

\subsection{Grupos de presión y poder político
(cont.)}\label{grupos-de-presiuxf3n-y-poder-poluxedtico-cont.-2}

\begin{itemize}
\tightlist
\item
  Entonces el nuevo votante pivotal (\emph{swing}) de cada grupo será:
\end{itemize}

\begin{equation}
\tilde{\sigma^{j}}=V^{j}(\mathbf{q_{A}})-V^{j}(\mathbf{q_{B}})-\tilde{\delta^{j}}+h(C_{B}-C_{A})
\end{equation}

\begin{itemize}
\tightlist
\item
  y encontramos que:
\end{itemize}

\begin{equation}
    p_{A}(\mathbf{q_{A}},\mathbf{q_{B}})=\frac{1}{2}+\psi
    \left[\sum_{j}
    \alpha^{j}[V^{j}(\mathbf{q_{A}})-V^{j}(\mathbf{q_{B}})]+h(C_{A}-C_{B})\right]
\end{equation}

\subsection{Grupos de presión y poder político
(cont.)}\label{grupos-de-presiuxf3n-y-poder-poluxedtico-cont.-3}

\begin{itemize}
\tightlist
\item
  Note la diferencia con modelos sin lobbies

  \begin{itemize}
  \tightlist
  \item
    no sólo debemos encontrar comportamiento óptimo de los candidatos
    sino también el comportamiento óptimo de los diferentes lobbies
  \end{itemize}
\item
  En cierta forma los candidatos anticipan los incentivos de los lobbies
  para darles recursos y eso debería reflejarse en las propuestas de
  política
\item
  Una funcion objetivo del lobbies puede ser maximizar utilidad esperada
  neta de contribuciones
\end{itemize}

\subsection{Grupos de presión y poder político
(cont.)}\label{grupos-de-presiuxf3n-y-poder-poluxedtico-cont.-4}

\begin{itemize}
\tightlist
\item
  Fn. objetivo de lobby \(j\) es:
\end{itemize}

\begin{equation}
p_{A}(\mathbf{q_{A}},\mathbf{q_{B}})V^{j}(\mathbf{q_{A}})+(1-p_{A}(\mathbf{q_{A}},\mathbf{q_{B}})V^{j}(\mathbf{q_{B}})-\frac{1}{2}\left((C^{j}_{B})^{2}+(C^{j}_{A})^{2}\right)
\end{equation}

\begin{itemize}
\tightlist
\item
  cada grupo busca maximizar la utilidad que le genera cada candidato
  ponderando por la probabilidad de que cada candidato sea elegido neta
  de costos asociados a la contribucion

  \begin{itemize}
  \tightlist
  \item
    ideología no juega; note que contribuciones ayudan a aumentar la
    probabilidad de ganar la eleccion
  \end{itemize}
\end{itemize}

\subsection{Grupos de presión y poder político
(cont.)}\label{grupos-de-presiuxf3n-y-poder-poluxedtico-cont.-5}

\begin{itemize}
\tightlist
\item
  La CPO respecto a la contribución para \(A\)
\end{itemize}

\begin{equation}
\frac{\partial{p_{A}}(\mathbf{q_{A}},\mathbf{q_{B}})}{\partial{C^{j}_{A}}}\left[V^{j}(\mathbf{q_{A}})-V^{j}(\mathbf{q_{B}})\right]=C^{j}_{A}
\end{equation}

\begin{itemize}
\tightlist
\item
  iguala el BMg de apoyar a \(A\) al CMg de proveer la contribución
  \(C^{j}_{A}\)
\end{itemize}

\begin{equation}
C^{j}_{A}=\max
\{0,h\psi\alpha^{j}[V^{j}(\mathbf{q_{A}})-V^{j}(\mathbf{q_{B}})]\}
\end{equation}

\subsection{Grupos de presión y poder político
(cont.)}\label{grupos-de-presiuxf3n-y-poder-poluxedtico-cont.-6}

\begin{quote}
\textbf{Problema de los candidatos.} Candidatos consideran incentivos de
\emph{lobbies} al fijar plataforma dado que esta afecta cuantas
contribuciones recien y a su vez esto incide sobre su probabilidad de
ser elegidos. El problema es perfectamente simétrico y ambos candidatos
convergen a la misma plataforma. Si no fuera así, un candidato puede
atraer \emph{toda} la contribución de algún lobby y el otro candidato
movería su plataforma y así.
\end{quote}

\subsection{Grupos de presión y poder político
(cont.)}\label{grupos-de-presiuxf3n-y-poder-poluxedtico-cont.-7}

\begin{itemize}
\tightlist
\item
  El candidato \(A\) elige la \(\mathbf{A}\) que maximiza:
\end{itemize}

\begin{align}
p_{A}(\mathbf{q_{A}},\mathbf{q_{B}})&=\frac{1}{2}+\psi
    \left[\sum_{j}\alpha^{j}[V^{j}(\mathbf{q_{A}})-V^{j}(\mathbf{q_{B}})]+h(\sum_{j}O^{j}\alpha^{j}C^{j}_{A}-\sum_{j}O^{j}\alpha^{j}C^{j}_{B})\right] \\
p_{A}(\mathbf{q_{A}},\mathbf{q_{B}})&=\frac{1}{2}+\psi
    \left[\sum_{j}\alpha^{j}[V^{j}(\mathbf{q_{A}})-V^{j}(\mathbf{q_{B}})]+h\sum_{j}O^{j}\alpha^{j}(C^{j}_{A}-C^{j}_{B})\right]
\end{align}

\subsection{Grupos de presión y poder político
(cont.)}\label{grupos-de-presiuxf3n-y-poder-poluxedtico-cont.-8}

\begin{itemize}
\tightlist
\item
  Sustituyendo las contribuciones óptimas de cada lobby:
\end{itemize}

\begin{align}
p_{A}(\mathbf{q_{A}},\mathbf{q_{B}})&=\frac{1}{2}+\psi
    \left[\sum_{j}\alpha^{j}[V^{j}(\mathbf{q_{A}})-V^{j}(\mathbf{q_{B}})]+h\sum_{j}O^{j}\alpha^{j}(h\psi\alpha^{j}[V^{j}(\mathbf{q_{A}})-V^{j}(\mathbf{q_{B}})])\right]
    \\
p_{A}(\mathbf{q_{A}},\mathbf{q_{B}})&=\frac{1}{2}+\psi
    \left[\sum_{j}\alpha^{j}[V^{j}(\mathbf{q_{A}})-V^{j}(\mathbf{q_{B}})]+h^{2}\psi\sum_{j}O^{j}\alpha^{j}(\alpha^{j}[V^{j}(\mathbf{q_{A}})-V^{j}(\mathbf{q_{B}})])\right]
\end{align}

\subsection{Grupos de presión y poder político
(cont.)}\label{grupos-de-presiuxf3n-y-poder-poluxedtico-cont.-9}

\begin{itemize}
\tightlist
\item
  Podemos tomar como dados todo lo que depende de \(\mathbf{q_{B}}\) y
  las constantes para maximizar:
\end{itemize}

\begin{align}
p_{A}(\mathbf{q_{A}},\mathbf{q_{B}})&=K+\psi
    \left[\sum_{j}\alpha^{j}[V^{j}(\mathbf{q_{A}})]+h^{2}\psi\sum_{j}O^{j}\alpha^{j}(\alpha^{j}[V^{j}(\mathbf{q_{A}})])\right] \\
p_{A}(\mathbf{q_{A}},\mathbf{q_{B}})&=K+\psi
    \left[\sum_{j}\alpha^{j}[V^{j}(\mathbf{q_{A}})][1+h^{2}\psi
    O^{j}\alpha^{j}]\right]
\end{align}

\subsection{Grupos de presión y poder político
(cont.)}\label{grupos-de-presiuxf3n-y-poder-poluxedtico-cont.-10}

\begin{quote}
La última ecuación muestra que el problema es equivalente a maximizar
una función de bienestar social ponderada en donde la ponderación de
cada grupo \(j\) viene dada por \(\alpha^{j}(1+h^{2}\psi
O^{j}\alpha^{j})\). Si todos los grupos tienen igual tamaño y todos
están organizados --\(\alpha^{j}=\alpha\) y \(O^{j}=1\), la solución
coincide con el optimo utilitarista. De lo contrario, grupos más grandes
recibirían mayor ponderación --i.e.~costo marginal de contribuciones es
creciente pero beneficio de contribuir es un bien público para todo el
grupo. Grupos más grandes reducen tamaño promedio de contribución por
miembro.
\end{quote}

\subsection{Grupos de presión y poder político
(cont.)}\label{grupos-de-presiuxf3n-y-poder-poluxedtico-cont.-11}

\begin{itemize}
\tightlist
\item
  Candidatos tienden a dar más peso a: 1) grupos más grandes; 2) con
  mayor organización

  \begin{itemize}
  \tightlist
  \item
    Si sólo los pobres se organizan, reciben mayor ponderación y
    ``extraen'' un mayor nivel de provisión de bien público --incluso
    mayor al óptimo (populismo no óptimo económicamente)
  \item
    Lo contrario si sólo los ricos se organizan
  \end{itemize}
\item
  ¿Qué grupos tendrán mayor probabilidad de organizarse? ¿Los pequeños?
  Posible menor problema de acción colectiva. ¿Los más grandes? ¿Los que
  más tienden a perder con una política pública? Beneficios
  concentrados, costos difusos.
\end{itemize}

\section{\texorpdfstring{\textbf{Un modelo de contribuciones de campaña
y \emph{lobbying} aplicado a datos de
Argentina}}{Un modelo de contribuciones de campaña y lobbying aplicado a datos de Argentina}}\label{un-modelo-de-contribuciones-de-campauxf1a-y-lobbying-aplicado-a-datos-de-argentina}

\section{\texorpdfstring{\textbf{Otros canales de influencia:
Corrupción}}{Otros canales de influencia: Corrupción}}\label{otros-canales-de-influencia-corrupciuxf3n}

\section{\texorpdfstring{\textbf{Conexiones políticas y personales.
Endorsements}}{Conexiones políticas y personales. Endorsements}}\label{conexiones-poluxedticas-y-personales.-endorsements}




\end{document}
